%!TEX root = ../GeneralRelativity.tex
% ──────────────────────────────────────────────────────────────
%  Addendum A — The energy-momentum tensor: a rigorous
%               distributional construction via the coarea formula
% ──────────────────────────────────────────────────────────────

\section{Addendum: the energy-momentum tensor from microscopics}%
\label{sec:emt-rigorous}

In Lecture~6 (\S\ref{sec:energy-momentum}) we defined the
energy-momentum tensor of a system of point particles via the
microscopic expressions~\eqref{eq:Tmunu-dt}
and~\eqref{eq:Tmunu-covariant}, and left the proof that these two
forms are equivalent---and that the result is a Lorentz tensor---as
an exercise (Exercise~\ref{ex:Tmunu-tensor}).  The standard argument,
due to Weinberg~\cite{Weinberg72}, consists of inserting a
coordinate-time integral, observing that ``the differentials~$dt'$
cancel,'' and replacing~$dt'$ with the invariant proper-time
element~$d\tau$.  While physically convincing, this argument glosses
over the distributional nature of the objects involved and leaves the
equivalence of the two forms essentially unproved.

In this addendum we supply the missing rigour.  The key tool turns
out to be the \textbf{coarea formula} of geometric measure
theory~\cite{Federer69}, which provides the natural framework for
slicing distributional objects supported on submanifolds by level
sets of smooth functions.  The one-dimensional specialisation of the
coarea formula is nothing but the classical change-of-variables
(substitution) formula, but recognising this connection places
Weinberg's construction on a firm mathematical footing and
generalises immediately to curved spacetimes and arbitrary Cauchy
surfaces.

% ──────────────────────────────────────────────────────────────
\subsection{Distributions supported on worldlines}%
\label{sec:dist-worldlines}

Throughout this addendum, $(\Rn{4}, \eta_{\alpha\beta})$ denotes
Minkowski spacetime with signature~$(-,+,+,+)$, and
$\mathcal{D}(\Rn{4}) = C_c^\infty(\Rn{4})$ is the space of
compactly supported smooth test functions.  We write
$\mathcal{D}'(\Rn{4})$ for the topological dual: the space of
distributions (continuous linear functionals on~$\mathcal{D}$).

\begin{definition}[Timelike worldline]\label{def:timelike-worldline}
  A \textbf{timelike worldline} is a smooth embedding
  $\gamma\colon I \to \Rn{4}$ (where $I\subseteq\R$ is an open
  interval) such that
  $\eta_{\alpha\beta}\,\dot\gamma^\alpha\dot\gamma^\beta < 0$
  everywhere.  We take $\tau$ to be the proper-time parameter and
  write $\dot\gamma^\alpha = d\gamma^\alpha/d\tau$.
\end{definition}

\noindent
For a massive particle of rest mass~$m$, the
\textbf{four-momentum} along~$\gamma$ is
$p^\alpha(\tau) = m\,\dot\gamma^\alpha(\tau)$.

The microscopic energy-momentum tensor is built from delta
functions localised on worldlines.  To make this precise, we define
it directly as a distribution.

\begin{definition}[Distributional energy-momentum tensor]%
\label{def:Tmunu-dist}
  Given $N$ timelike worldlines $\gamma_j\colon \R \to \Rn{4}$
  carrying four-momenta $p_j^\alpha(\tau)$, the
  \textbf{(distributional) energy-momentum tensor}
  $T^{\alpha\beta} \in \mathcal{D}'(\Rn{4})$ is defined by its
  action on test functions:
  \begin{equation}\label{eq:Tmunu-dist-def}
    \eqbox{\langle T^{\alpha\beta},\, \phi\rangle
      \;=\; \sum_{j=1}^{N} \int_{\R} p_j^\alpha(\tau)\,
        \dot\gamma_j^\beta(\tau)\,
        \phi\bigl(\gamma_j(\tau)\bigr)\;\dd\tau}
  \end{equation}
  for every $\phi \in \mathcal{D}(\Rn{4})$.
\end{definition}

\begin{proposition}[Well-definedness]\label{prop:Tmunu-well-defined}
  The functional~\eqref{eq:Tmunu-dist-def} defines a distribution
  of order zero, that is, a tensor-valued Radon measure
  on~$\Rn{4}$.
\end{proposition}

\begin{proof}
  For each~$j$, the worldline~$\gamma_j$ is timelike, so the
  time component $\gamma_j^0(\tau)$ satisfies
  $d\gamma_j^0/d\tau > 0$ (future-directed).  Hence the map
  $\tau \mapsto \gamma_j^0(\tau)$ is a diffeomorphism
  from~$I$ onto its image, and $\gamma_j$ is a proper map: the
  preimage $\gamma_j^{-1}(K)$ of any compact set
  $K \subset \Rn{4}$ is compact.

  Now let $\phi \in \mathcal{D}(\Rn{4})$ have
  support in~$K$.  The function
  $\tau \mapsto p_j^\alpha(\tau)\,\dot\gamma_j^\beta(\tau)\,
  \phi(\gamma_j(\tau))$ is smooth with support in the compact set
  $\gamma_j^{-1}(K)$, so the integral converges absolutely.
  Moreover,
  \[
    |\langle T^{\alpha\beta}, \phi\rangle|
      \;\leq\; \sum_j \int_{\gamma_j^{-1}(K)}
        |p_j^\alpha|\,|\dot\gamma_j^\beta|\;\dd\tau
        \cdot \|\phi\|_\infty
      \;\leq\; C_K\,\|\phi\|_\infty\,,
  \]
  so $T^{\alpha\beta}$ is continuous in the $C^0$-topology and
  hence is a distribution of order zero.
\end{proof}

\begin{remark}
  Definition~\ref{def:Tmunu-dist} is the \emph{pushforward} of a
  measure on the worldline to a distribution on spacetime.  In the
  language of de~Rham's theory of currents, it defines a
  $1$-current on~$\Rn{4}$---the integration current of
  the worldline, weighted by the tensor
  $p^\alpha\,\dd\gamma^\beta$.
  See Federer~\cite[\S4.1]{Federer69}.
\end{remark}

\noindent
Informally, Definition~\ref{def:Tmunu-dist} says that
$T^{\alpha\beta}$ is the distribution
\begin{equation}\label{eq:Tmunu-delta4-add}
  T^{\alpha\beta}(x) = \sum_{j=1}^{N} \int\!\dd\tau\;
    p_j^\alpha(\tau)\,\dot\gamma_j^\beta(\tau)\,
    \delta^{(4)}\!\bigl(x - \gamma_j(\tau)\bigr)\,,
\end{equation}
which is precisely equation~\eqref{eq:Tmunu-covariant}.  The
distributional definition makes it clear that this expression is
well defined and coordinate-independent; what remains is to show
that it reduces to the $3{+}1$
form~\eqref{eq:Tmunu-dt} upon slicing by a constant-time
hyperplane.

% ──────────────────────────────────────────────────────────────
\subsection{The coarea formula}\label{sec:coarea}

The bridge between the covariant
form~\eqref{eq:Tmunu-delta4-add} and the $3{+}1$
splitting~\eqref{eq:Tmunu-dt} is provided by the
\textbf{coarea formula}.  We state the general result and then
specialise to the one-dimensional case that we need.

\begin{theorem}[Coarea formula {\cite[Theorem~3.2.12]{Federer69}}]%
\label{thm:coarea}
  Let $f\colon \Rn{n} \to \Rn{m}$ be a Lipschitz map with
  $m \leq n$.  Then for any nonnegative Borel function
  $g\colon \Rn{n} \to [0,\infty]$,
  \begin{equation}\label{eq:coarea-general}
    \int_{\Rn{n}} g(x)\,J_m f(x)\;\dd\mathcal{H}^n(x)
      = \int_{\Rn{m}}
        \biggl(\int_{f^{-1}(y)}
          g\;\dd\mathcal{H}^{n-m}\biggr)
        \dd\mathcal{H}^m(y)\,,
  \end{equation}
  where $J_m f$ denotes the $m$-dimensional Jacobian of~$f$
  and $\mathcal{H}^k$ the $k$-dimensional Hausdorff measure.
\end{theorem}

\noindent
When $n = m = 1$, the Hausdorff measures reduce to Lebesgue
measure, the $1$-dimensional Jacobian is simply $|f'|$, and the
level set $f^{-1}(t)$ is discrete.  The coarea formula becomes:

\begin{corollary}
  Let $f\colon \R \to \R$ be Lipschitz and
  $g\colon \R \to [0,\infty]$ Borel.  Then
  \begin{equation}\label{eq:coarea-1d}
    \eqbox{
    \int_{\R} g(\tau)\,|f'(\tau)|\;\dd\tau
      = \int_{\R}
        \sum_{\tau\in f^{-1}(t)} g(\tau)\;\dd t}\,.
  \end{equation}
\end{corollary}

\begin{remark}\label{rem:delta-composition}
  When $f$ is a diffeomorphism---as will be the case for the time
  coordinate along a timelike worldline---each level set
  $f^{-1}(t)$ contains exactly one point~$\tau_\ast(t)$, and
  \eqref{eq:coarea-1d} reduces to the elementary substitution
  $g(\tau)\,|f'(\tau)|\,\dd\tau = g(\tau_\ast(t))\,\dd t$.
  This is equivalent to the distributional identity
  \begin{equation}\label{eq:delta-composition}
    \delta\bigl(f(\tau) - t\bigr)
      = \frac{\delta(\tau - \tau_\ast)}{|f'(\tau_\ast)|}\,,
  \end{equation}
  so the coarea formula provides the rigorous foundation for the
  ``delta function composition rule'' that is often invoked without
  proof in physics texts.
\end{remark}

% ──────────────────────────────────────────────────────────────
\subsection{From the covariant form to the spatial-slice form}%
\label{sec:equivalence}

We now prove the main result: the covariant
definition~\eqref{eq:Tmunu-delta4-add} reproduces the $3{+}1$
form~\eqref{eq:Tmunu-dt} of Lecture~6 upon slicing by a
constant-time hyperplane.

\begin{theorem}[Equivalence of forms]\label{thm:equivalence}
  Let $T^{\alpha\beta}$ be the distribution defined
  in~\eqref{eq:Tmunu-dist-def}.  Then, in the sense of
  distributions on~$\Rn{4}$,
  \begin{equation}\label{eq:equivalence}
    T^{\alpha\beta}(x) = \sum_{j=1}^{N} p_j^\alpha(t)\,
      \frac{d\gamma_j^\beta}{dt}(t)\,
      \delta^{(3)}\!\bigl(\vb{x} - \boldsymbol{\gamma}_j(t)\bigr)\,,
  \end{equation}
  where $t = x^0$ and $\boldsymbol{\gamma}_j(t) =
  \bigl(\gamma_j^1(t),\gamma_j^2(t),\gamma_j^3(t)\bigr)$ is the
  spatial position of the $j$-th particle at coordinate time~$t$.
\end{theorem}

\begin{proof}
  It suffices to prove the identity for a single worldline; the
  sum over~$j$ follows by linearity.  Drop the subscript~$j$ and
  let $\gamma\colon \R \to \Rn{4}$ be a timelike worldline
  carrying momentum~$p^\alpha$.

  \medskip\noindent
  \textbf{Step~1: Factor out the time velocity.}\;
  Define the \emph{time function along the worldline}
  $f\colon \R \to \R$ by $f(\tau) = \gamma^0(\tau)$.
  Since $\gamma$ is future-directed and timelike,
  $f'(\tau) = d\gamma^0/d\tau > 0$, so $f$ is a smooth
  diffeomorphism from~$\R$ onto~$\R$.

  Decompose the four-velocity:
  \[
    \dot\gamma^\beta(\tau)
      = \frac{d\gamma^\beta}{d\tau}
      = \frac{d\gamma^\beta}{dt}\,\frac{dt}{d\tau}
      = \frac{d\gamma^\beta}{dt}\,f'(\tau)\,.
  \]
  Substituting into Definition~\ref{def:Tmunu-dist}:
  \begin{equation}\label{eq:step1}
    \langle T^{\alpha\beta},\phi\rangle
      = \int_{\R} p^\alpha(\tau)\,
        \frac{d\gamma^\beta}{dt}(\tau)\,
        \phi\bigl(\gamma(\tau)\bigr)\,f'(\tau)\;\dd\tau\,.
  \end{equation}

  \medskip\noindent
  \textbf{Step~2: Apply the coarea formula.}\;
  The integrand in~\eqref{eq:step1} has the form
  $g(\tau)\,|f'(\tau)|\;\dd\tau$ with $|f'| = f'$ (since
  $f' > 0$).  By the one-dimensional coarea
  formula~\eqref{eq:coarea-1d}, the level set
  $f^{-1}(t) = \{\tau_\ast(t)\}$ is a singleton for each~$t$, and
  \begin{equation}\label{eq:step2}
    \langle T^{\alpha\beta},\phi\rangle
      = \int_{\R} p^\alpha\bigl(\tau_\ast(t)\bigr)\,
        \frac{d\gamma^\beta}{dt}\bigl(\tau_\ast(t)\bigr)\,
        \phi\bigl(\gamma(\tau_\ast(t))\bigr)\;\dd t\,.
  \end{equation}

  \medskip\noindent
  \textbf{Step~3: Identify the spatial delta function.}\;
  By definition of $\tau_\ast(t)$, we have
  $\gamma^0(\tau_\ast(t)) = t$, so
  $\gamma(\tau_\ast(t)) = \bigl(t,\,\boldsymbol{\gamma}(t)\bigr)$.
  Writing $p^\alpha(\tau_\ast(t)) = p^\alpha(t)$ and
  $d\gamma^\beta/dt\,\big|_{\tau_\ast(t)} = d\gamma^\beta/dt\,(t)$
  as functions of coordinate time, equation~\eqref{eq:step2}
  becomes
  \begin{align}
    \langle T^{\alpha\beta},\phi\rangle
      &= \int_{\R} p^\alpha(t)\,
         \frac{d\gamma^\beta}{dt}(t)\,
         \phi\bigl(t,\boldsymbol{\gamma}(t)\bigr)\;\dd t
         \notag\\
      &= \int_{\R}\dd t\int_{\Rn{3}}\dd^3\!\vb{x}\;\;
         p^\alpha(t)\,\frac{d\gamma^\beta}{dt}(t)\,
         \delta^{(3)}\!\bigl(\vb{x} - \boldsymbol{\gamma}(t)\bigr)\,
         \phi(t,\vb{x})\,.
         \label{eq:step3}
  \end{align}
  Since~\eqref{eq:step3} holds for all
  $\phi\in\mathcal{D}(\Rn{4})$, we conclude
  that~\eqref{eq:equivalence} holds in the distributional sense.
\end{proof}

\begin{intuition}[Slicing worldlines by time]
  The theorem says that the natural (covariant) definition of
  $T^{\alpha\beta}$---as a measure spread along the worldline---can
  be ``sliced'' by a constant-time hyperplane $\Sigma_t = \{x^0 = t\}$
  to recover the $3{+}1$ expression.  The coarea formula provides
  exactly the right Jacobian factor ($1/f'(\tau_\ast)$) to convert
  the proper-time line element~$\dd\tau$ along the worldline into
  the transverse coordinate~$\dd t$.  This factor cancels the
  $d\gamma^\beta/d\tau$ in the covariant integrand, producing
  $d\gamma^\beta/dt$ in the sliced expression.

  Weinberg's observation that ``the $dt'$ cancel'' is a compact
  (if informal) summary of exactly this mechanism.
\end{intuition}

% ──────────────────────────────────────────────────────────────
\subsection{Lorentz covariance}\label{sec:lorentz-covariance-add}

The covariant form~\eqref{eq:Tmunu-dist-def} immediately yields
the transformation law, resolving
Exercise~\ref{ex:Tmunu-tensor}:

\begin{proposition}[Lorentz covariance]\label{prop:lorentz-cov}
  Under a Lorentz transformation
  $x'^\mu = \Lambda^\mu{}_\nu\, x^\nu$, the distributional
  energy-momentum tensor transforms as
  \begin{equation}\label{eq:Tmunu-lorentz}
    T'^{\alpha\beta}(x')
      = \Lambda^\alpha{}_\gamma\,
        \Lambda^\beta{}_\delta\,
        T^{\gamma\delta}(x)\,.
  \end{equation}
\end{proposition}

\begin{proof}
  In the covariant form~\eqref{eq:Tmunu-dist-def}, the integrand
  consists of:
  \begin{itemize}
    \item $p^\alpha(\tau)$: a four-vector,
    \item $\dot\gamma^\beta(\tau) = d\gamma^\beta/d\tau$: a
      four-vector,
    \item $\phi(\gamma(\tau))$: a scalar,
    \item $\dd\tau$: Lorentz-invariant.
  \end{itemize}
  Thus $\langle T^{\alpha\beta}, \phi\rangle$ carries two
  contravariant Lorentz indices, and the transformation
  law~\eqref{eq:Tmunu-lorentz} follows.  (One also uses
  $|\det\Lambda| = 1$, which ensures that the four-dimensional
  delta function $\delta^{(4)}$ is a Lorentz scalar.)
\end{proof}

\begin{remark}
  The $3{+}1$ form~\eqref{eq:Tmunu-dt} is \emph{not} manifestly
  covariant: neither $\delta^{(3)}$ nor $d\gamma^\beta/dt$ is
  Lorentz-invariant in isolation.  It is only through
  Theorem~\ref{thm:equivalence}---whose proof rests on the coarea
  formula---that we can deduce that the combination
  $p^\alpha\,(d\gamma^\beta/dt)\,\delta^{(3)}$ transforms as a
  tensor.  The miracle is that the Jacobian factors conspire to
  produce a Lorentz-invariant measure along the worldline; the
  coarea formula explains \emph{why} they conspire.
\end{remark}

% ──────────────────────────────────────────────────────────────
\subsection{Distributional conservation law}\label{sec:conservation-add}

We now prove the conservation law entirely within the
distributional framework, avoiding the somewhat formal manipulation
of derivatives of delta functions that appears in
Weinberg~\cite[\S2.8]{Weinberg72}.

\begin{theorem}[Conservation of $T^{\alpha\beta}$]%
\label{thm:conservation}
  If each particle is free ($\dd p_j^\alpha/\dd\tau = 0$), then
  \begin{equation}\label{eq:conservation}
    \partial_\beta\, T^{\alpha\beta} = 0
  \end{equation}
  in the distributional sense.  More generally, in the presence of
  external forces,
  \begin{equation}\label{eq:conservation-force}
    \partial_\beta\, T^{\alpha\beta}(x) = G^\alpha(x)\,,
  \end{equation}
  where the \textbf{force density}
  \[
    G^\alpha(x) = \sum_j \int\!\dd\tau\;
      \frac{\dd p_j^\alpha}{\dd\tau}(\tau)\,
      \delta^{(4)}\!\bigl(x - \gamma_j(\tau)\bigr)
  \]
  is the distributional pushforward of the equations of motion.
\end{theorem}

\begin{proof}
  We prove~\eqref{eq:conservation};
  equation~\eqref{eq:conservation-force} follows by the same
  argument without discarding the boundary term.  For any
  $\phi \in \mathcal{D}(\Rn{4})$:
  \begin{align}
    \langle \partial_\beta T^{\alpha\beta},\,\phi\rangle
      &= -\langle T^{\alpha\beta},\,\partial_\beta\phi\rangle
      \notag\\
      &= -\sum_j \int_{\R} p_j^\alpha(\tau)\,
        \dot\gamma_j^\beta(\tau)\,
        \partial_\beta\phi\bigl(\gamma_j(\tau)\bigr)\;\dd\tau
      \notag\\
      &= -\sum_j \int_{\R} p_j^\alpha(\tau)\,
        \frac{\dd}{\dd\tau}\bigl[\phi(\gamma_j(\tau))\bigr]
        \;\dd\tau\,,
        \label{eq:cons-chain}
  \end{align}
  where the last step uses the chain rule
  $\frac{\dd}{\dd\tau}\phi(\gamma(\tau))
  = \dot\gamma^\beta(\tau)\,\partial_\beta\phi(\gamma(\tau))$.
  Integrating by parts (the boundary terms vanish since $\phi$
  has compact support and the worldline eventually leaves any
  compact set):
  \begin{equation}\label{eq:cons-ibp}
    \langle \partial_\beta T^{\alpha\beta},\,\phi\rangle
      = \sum_j \int_{\R}
        \frac{\dd p_j^\alpha}{\dd\tau}(\tau)\,
        \phi\bigl(\gamma_j(\tau)\bigr)\;\dd\tau\,.
  \end{equation}
  For free particles, $\dd p_j^\alpha/\dd\tau = 0$, and the
  right-hand side vanishes for all~$\phi$.
\end{proof}

\begin{remark}
  The elegance of the distributional proof is worth
  emphasising.  In the $3{+}1$ approach, the conservation law
  requires differentiating $\delta^{(3)}\!\bigl(\vb{x} -
  \boldsymbol{\gamma}(t)\bigr)$ with respect to~$\vb{x}$ and
  then using the chain rule
  $\frac{\partial}{\partial x^i}\,\delta^{(3)}
  = -\frac{\partial}{\partial \gamma^i}\,\delta^{(3)}
  = -\frac{1}{\dot\gamma^i}\,\frac{\partial}{\partial t}\,
  \delta^{(3)}$
  to convert spatial derivatives into time derivatives.  The
  distributional proof~\eqref{eq:cons-chain}--\eqref{eq:cons-ibp}
  bypasses this entirely: the chain rule on test functions and
  integration by parts do all the work.
\end{remark}

\begin{exercise}\label{ex:conservation-collision}
  Extend Theorem~\ref{thm:conservation} to particles that undergo
  strictly localised collisions (so that worldlines begin and end
  at discrete spacetime events).  Show that $\partial_\beta
  T^{\alpha\beta} = 0$ provided that four-momentum is conserved
  at each collision vertex: $\sum_{n\in c} p_n^\alpha = 0$ for
  each collision~$c$.
  \emph{Hint:} integrate by parts on each worldline segment;
  the boundary terms at each collision vertex sum to zero by
  momentum conservation.
\end{exercise}

% ──────────────────────────────────────────────────────────────
\subsection{Slicing by an arbitrary spacelike hypersurface}%
\label{sec:general-slicing}

The equivalence of the covariant and $3{+}1$ forms does not depend
on the particular choice of time coordinate~$x^0$.  We now state
the general result, which applies to arbitrary spacelike
hypersurfaces and extends naturally to curved spacetimes.

\begin{theorem}[General slicing theorem]%
\label{thm:general-slicing}
  Let $\Sigma = f^{-1}(0)$ be a smooth spacelike hypersurface
  in~$\Rn{4}$ defined as a level set of a smooth function
  $f\colon\Rn{4}\to\R$ with $\nabla f$ everywhere timelike.
  Then each worldline~$\gamma_j$ intersects~$\Sigma$
  transversally at a unique point
  $\gamma_j(\tau_j^\ast)$, and the restriction (``slice'') of
  $T^{\alpha\beta}$ to~$\Sigma$ satisfies
  \begin{equation}\label{eq:general-slicing}
    T^{\alpha\beta}\big|_\Sigma(\vb{y})
      = \sum_{j=1}^{N}
        \frac{p_j^\alpha(\tau_j^\ast)\,
              \dot\gamma_j^\beta(\tau_j^\ast)}
             {\bigl|\partial_\gamma f\,
              \dot\gamma_j^\gamma(\tau_j^\ast)\bigr|}\,
        \delta_\Sigma(\vb{y} - \vb{y}_j)\,,
  \end{equation}
  where $\vb{y}$ denotes coordinates on~$\Sigma$, $\vb{y}_j$ is
  the location of the $j$-th crossing on~$\Sigma$, and
  $\delta_\Sigma$ is the delta function with respect to the induced
  measure on~$\Sigma$.
\end{theorem}

\begin{proof}[Sketch of proof]
  Apply the slicing theorem for normal $1$-currents
  (Federer~\cite[Theorem~4.3.2]{Federer69}) to the integration
  current associated with each worldline, with slicing
  function~$f$.  The transversality condition
  $\partial_\gamma f \cdot \dot\gamma^\gamma \neq 0$ is guaranteed
  because $\nabla f$ is timelike and $\dot\gamma$ is timelike,
  so their Lorentzian inner product is nonzero by the reverse
  Cauchy--Schwarz inequality.  The Jacobian factor
  $|\partial_\gamma f \cdot \dot\gamma^\gamma|^{-1}$ arises from
  the coarea formula applied to the composition
  $f\circ\gamma\colon\R\to\R$.
\end{proof}

\begin{remark}
  For the standard time function $f(x) = x^0 - t_0$, we have
  $\partial_\gamma f = \delta_\gamma^0$, so
  $|\partial_\gamma f\cdot\dot\gamma^\gamma| = |\dot\gamma^0|
  = d\gamma^0/d\tau$, and the Jacobian factor combines with
  $\dot\gamma^\beta$ to give $d\gamma^\beta/dt$---recovering
  Theorem~\ref{thm:equivalence}.  In curved spacetime one
  replaces $f$ with a Cauchy time function and $\delta_\Sigma$
  with the delta function defined by the induced Riemannian metric
  on the Cauchy surface.
\end{remark}

% ──────────────────────────────────────────────────────────────
\subsection{The four-current revisited}\label{sec:four-current-add}

For completeness, we note that the entire development applies
\emph{mutatis mutandis} to the electromagnetic four-current
$J^\alpha$ of \S\ref{sec:four-current}.  Given particles with
charges~$q_j$, define the distributional four-current by
\begin{equation}\label{eq:Jmu-dist}
  \langle J^\alpha,\,\phi\rangle
    = \sum_{j=1}^{N} \int_{\R} q_j\,
      \dot\gamma_j^\alpha(\tau)\,
      \phi\bigl(\gamma_j(\tau)\bigr)\;\dd\tau\,.
\end{equation}
The coarea argument of Theorem~\ref{thm:equivalence} shows that
\begin{equation}\label{eq:Jmu-31}
  J^\alpha(x) = \sum_{j=1}^{N} q_j\,
    \frac{d\gamma_j^\alpha}{dt}(t)\,
    \delta^{(3)}\!\bigl(\vb{x} - \boldsymbol{\gamma}_j(t)\bigr)\,,
\end{equation}
which unifies the charge density~\eqref{eq:charge-density} and
current density~\eqref{eq:current-density}.  The conservation law
$\partial_\alpha J^\alpha = 0$ follows from the same
integration-by-parts argument as
Theorem~\ref{thm:conservation}, using the constancy of charge
along each worldline ($\dd q_j/\dd\tau = 0$).

% ──────────────────────────────────────────────────────────────
\subsection{Discussion}\label{sec:emt-discussion}

\begin{keyresult}[Summary]
  \begin{enumerate}
    \item[\textup{(i)}] The covariant
      form~\eqref{eq:Tmunu-covariant} is the fundamental
      definition of the microscopic energy-momentum tensor: it is a
      well-defined distribution
      (Proposition~\ref{prop:Tmunu-well-defined}) and a Lorentz
      tensor by construction
      (Proposition~\ref{prop:lorentz-cov}).
    \item[\textup{(ii)}] The $3{+}1$ form~\eqref{eq:Tmunu-dt} is
      derived from the covariant form by \emph{slicing with the
      coarea formula} (Theorem~\ref{thm:equivalence}).
    \item[\textup{(iii)}] Conservation
      $\partial_\beta T^{\alpha\beta} = 0$ follows from the
      equations of motion by a distributional
      integration-by-parts argument
      (Theorem~\ref{thm:conservation}), without manipulating
      derivatives of delta functions.
    \item[\textup{(iv)}] The construction generalises immediately
      to arbitrary spacelike hypersurfaces via the slicing
      theorem for currents
      (Theorem~\ref{thm:general-slicing}).
  \end{enumerate}
\end{keyresult}

\noindent
Thus Weinberg's beautiful microscopic picture of the
energy-momentum tensor---as the density and flux of four-momentum
carried by point particles---rests on a rigorous distributional
foundation.  The coarea formula, far from being an exotic tool of
geometric measure theory, is the elementary mechanism (change of
variables along the worldline) that makes the entire construction
work.  Recognising this clarifies what would otherwise appear to
be a sleight of hand in the passage from coordinate-time
expressions to manifestly covariant ones.
